\documentclass[11pt]{article}
\usepackage[margin=1in]{geometry}
\usepackage{amsmath,amssymb,amsthm}
\usepackage{booktabs}
\usepackage{hyperref}
\hypersetup{colorlinks=true,linkcolor=blue,citecolor=blue,urlcolor=blue}

\newtheorem{theorem}{Theorem}
\newtheorem{proposition}[theorem]{Proposition}
\newtheorem{lemma}[theorem]{Lemma}
\newtheorem{corollary}[theorem]{Corollary}
\newtheorem{conjecture}{Conjecture}
\renewcommand{\theconjecture}{\Alph{conjecture}}
\theoremstyle{definition}
\newtheorem{definition}[theorem]{Definition}
\theoremstyle{remark}
\newtheorem{remark}[theorem]{Remark}

\newcommand{\Pp}{\mathbb{P}}

\title{Goldbach Partitions with a Twin, Cousin or Sexy Starter:\\
a Mod-3 Obstruction, a Density Law, and Verification to $10^{13}$}

\author{Santhiagu Antony Raj Parthibanathan\\
\small Independent researcher, Tamil Nadu, India\\
\small ORCID: \href{https://orcid.org/0009-0002-0355-4191}{0009-0002-0355-4191}
\quad Email: \href{mailto:psaraj12@gmail.com}{\texttt{psaraj12@gmail.com}}}

\date{\today}

\begin{document}
\maketitle

\begin{abstract}
Call an odd prime $p$ a \emph{twin starter} if $p+2$ is prime, a \emph{cousin
starter} if $p+4$ is prime, and a \emph{sexy starter} if $p+6$ is prime. We
study two strengthenings of the binary Goldbach conjecture: (A) every even
$N\ge 6$ is $p+q$ with $q$ prime and $p$ a twin or cousin starter; (B) every
even $N\ge 8$ is $p+q$ with $q$ prime and $p$ a sexy starter. Both are verified
exhaustively with a dedicated C++/OpenMP verifier: (A) for all even
$6\le N\le 10^{13}$ and (B) for all even $8\le N\le 10^{12}$, with every witness
satisfying $p\le 2\cdot 10^7$; neither has a counterexample. We explain why
(A) needs two gaps while (B) needs one. Twin starters alone fail
\emph{exactly} for the even $N\equiv 2\pmod 6$ with $N-3$ composite, and this
set has relative density $\tfrac13-(1+o(1))/\ln X$ among even $N\le X$, an
unconditional consequence of the prime number theorem for arithmetic
progressions. Cousin starters alone fail exactly on the mirror class. The
prime $3$ is the only modulus that can block a whole residue class in this
way, and it does so precisely when $3\nmid g$ for the gap $g$; for $g=6$ it
only halves the usable starters. The measured fraction of even $N$ lacking a
twin-starter partition agrees with the predicted law
$\tfrac13\,(1-3/\ln N)$ to within sampling error from $10^4$ to $10^{13}$.
\end{abstract}

\section{Introduction}

\paragraph{The idea in plain words.}
Goldbach's conjecture says every even number from $4$ upward is a sum of two
primes. Look at the partitions of a large even number and a striking pattern
appears: one of the two primes can almost always be chosen \emph{very small},
and that small prime can almost always be chosen to sit right next to another
prime. For example, it might be a $p$ with $p+2$ also prime. It is as if the
partitions of consecutive even numbers were pushed forward by a small engine
built out of nearby prime pairs: a partition of $N$ with $q$ prime and $q+2$
prime immediately gives a partition of $N+2$. This paper turns that picture
into two precise conjectures, verifies them over trillions of cases, and
explains exactly which prime pairs the engine needs and why. The answer
comes down to arithmetic modulo $3$.

\paragraph{Statements.}

\begin{definition}
An odd prime $p$ is a \emph{twin starter} if $p+2\in\Pp$, a \emph{cousin
starter} if $p+4\in\Pp$, and a \emph{sexy starter} if $p+6\in\Pp$. More
generally, for even $g\ge 2$ it is a \emph{$g$-starter} if $p+g\in\Pp$.
\end{definition}

\begin{conjecture}[twin/cousin starters]\label{conj:A}
Every even $N\ge 6$ can be written $N=p+q$ with $q$ prime and $p$ a twin
starter or a cousin starter.
\end{conjecture}

\begin{conjecture}[sexy starters]\label{conj:B}
Every even $N\ge 8$ can be written $N=p+q$ with $q$ prime and $p$ a sexy
starter.
\end{conjecture}

Both conjectures imply Goldbach's conjecture for the corresponding range. Both
also imply that the starter sets are infinite: if only finitely many starters
existed, $N-p$ would eventually be forced to take a composite value for some $N$.
(For instance, by choosing $N$ with $N-p$ divisible by a fixed prime for every
one of the finitely many $p$; see Remark~\ref{rem:infinitude}.) So
Conjecture~\ref{conj:A} implies there are infinitely many twin or cousin
prime pairs, and Conjecture~\ref{conj:B} implies infinitely many sexy prime
pairs. Neither statement is proved.

In every verified case the witness $p$ is tiny compared with $N$. In our sampled
windows up to $10^{13}$ the least starter partner never exceeded $6{,}197$; exhaustively
it is larger, at least $20{,}747$ below $10^{13}$ (Section~\ref{sec:fuel}), but still far
below the anchor limit. In particular the witness is always the
\emph{smaller} part of the partition for $N>4\cdot 10^7$.

\paragraph{Related work.}
Representations of even numbers using twin primes go back at least to
Zwillinger~\cite{Zwillinger}. Dubner~\cite{Dubner} conjectured that every even
$n>4208$ is a sum of two twin-prime \emph{members} (OEIS
A007534~\cite{A007534}). Sahu~\cite{Sahu} and, independently and much more
extensively, Pandey~\cite{Pandey} studied the intermediate statement that every
even $n\ge 6$ is $p+t$ with $p$ prime and $t$ a member (lower or upper) of a
twin pair. Pandey verified it to $10^{14}$, observed the same mod-3
``orientation rigidity'' that underlies Section~\ref{sec:mod3}, proved by Romanov's method
that it holds for a positive proportion of even numbers under a Hardy--Littlewood-type
lower bound for twin primes, and modelled the size of the least witness. Gómez~\cite{Gomez}
proposed weaker twin, cousin and sexy variants (either summand may be either
member of its pair), checked up to about $6.5\cdot 10^9$. Sun~\cite{Sun}
formulated many refinements of Goldbach involving prime pairs, including sexy
pairs, though with different side conditions. Oliveira e Silva, Herzog and
Pardi~\cite{OHP} verified the binary Goldbach conjecture to $4\cdot 10^{18}$, and the
author extended this to $4.003\cdot10^{18}$ with the Q-Hot method~\cite{QHot2}, whose
verifier is reused here (Section~\ref{sec:computation}). The forward-propagation
step of the introduction (if $N=p+q$ and $p+2k$ is prime, then
$N+2k=(p+2k)+q$) is the basis of that method and was first described in the
author's April 2026 preprint~\cite{QHot}.

As far as we can tell, the precise Conjectures~\ref{conj:A} and~\ref{conj:B}
(lower members of gap-$2$/gap-$4$ pairs, and lower members of gap-$6$ pairs,
respectively), the density law of Theorem~\ref{thm:law}, and the
classification of gaps by $g \bmod 3$ (Proposition~\ref{prop:only3}) have not
appeared before. We welcome corrections. The new contribution relative
to~\cite{Pandey} is that Conjecture~\ref{conj:A} covers the blocked residue
class with \emph{cousin} pairs rather than upper twin members, together with the
single-gap Conjecture~\ref{conj:B} and the quantitative results below.

\section{The mod-3 obstruction}\label{sec:mod3}

\begin{lemma}\label{lem:residues}
Let $p>3$ be prime.
\begin{enumerate}
\item If $p$ is a twin starter then $p\equiv 2\pmod 3$.
\item If $p$ is a cousin starter then $p\equiv 1\pmod 3$.
\item A sexy starter may lie in either nonzero class modulo $3$.
\item $3$ is both a twin starter $(3,5)$ and a cousin starter $(3,7)$, and it is
the only prime that is both. It is not a sexy starter, since $9$ is composite.
\end{enumerate}
\end{lemma}

\begin{proof}
If $p\equiv1$ then $p+2\equiv0\pmod 3$, so $p+2$ is not prime unless it equals $3$; this
proves (1), and (2) is identical with $p+4$. For (3), $p+6\equiv p\pmod 3$.
For (4), among $p,p+2,p+4$ one is divisible by $3$, so all three are prime only
for $p=3$.
\end{proof}

\begin{proposition}\label{prop:twinblock}
Let $N\ge 8$ be even.
\begin{enumerate}
\item If $N\equiv 2\pmod 6$, then $N=p+q$ with $q$ prime and $p$ a twin
starter if and only if $N-3$ is prime.
\item If $N\equiv 4\pmod 6$, then $N=p+q$ with $q$ prime and $p$ a cousin
starter if and only if $N-3$ is prime.
\end{enumerate}
\end{proposition}

\begin{proof}
(1) If $N-3$ is prime, take $p=3$. Conversely, let $N=p+q$ with $p$ a twin
starter and $q$ prime. If $p=3$ then $q=N-3$. If $p>3$ then
$p\equiv 2\pmod 3$ by Lemma~\ref{lem:residues}, so $q=N-p\equiv 0\pmod 3$, hence
$q=3$ and $p=N-3$ is prime. (2) is identical, using $p\equiv1\pmod3$ for
cousin starters $p>3$.
\end{proof}

So neither gap works alone. Twin starters fail for every $N\equiv 2\pmod 6$ with
$N-3$ composite, and cousin starters fail for every $N\equiv 4\pmod 6$ with
$N-3$ composite. Both sets are infinite and, by Theorem~\ref{thm:law}, each
has relative density $1/3$ among even numbers. In the other direction,
Proposition~\ref{prop:twinblock} turns Conjecture~\ref{conj:A} into three
cleaner statements, one per residue class. We call this the \emph{gearbox}.

\begin{corollary}[gearbox form]\label{cor:gearbox}
Conjecture~\ref{conj:A} is equivalent to the conjunction of:
\begin{enumerate}
\item[(a)] every even $N\equiv 0\pmod 6$ has a twin-starter or cousin-starter
partition;
\item[(b)] every even $N\equiv 4\pmod 6$ has a twin-starter partition;
\item[(c)] every even $N\equiv 2\pmod 6$, $N\ge 8$, has a cousin-starter partition.
\end{enumerate}
\end{corollary}

\begin{proof}
For $N\equiv 4\pmod 6$, a cousin-starter partition exists only when $N-3$ is
prime (Proposition~\ref{prop:twinblock}(2)), and then $p=3$ is also a twin
starter. So in this class Conjecture~\ref{conj:A} is equivalent to (b). The
class $N\equiv 2\pmod 6$ is symmetric.
\end{proof}

Empirically the gearbox is rigid. Up to $10^8$, the even $N$ with no
twin-starter partition are \emph{exactly} the $N\equiv 2\pmod 6$ with $N-3$
composite (exhaustive check), and the same holds for every one of the
$1.4\times10^6$ sampled $N$ up to $10^{13}$ (Section~\ref{sec:fuel}).

\begin{proposition}[only $3$ obstructs]\label{prop:only3}
Let $g\ge 2$ be even and let $\ell$ be a prime. Call a residue class
$N\bmod \ell$ \emph{blocked for gap $g$} if every $g$-starter $p$ with $p\ne\ell$
and $p+g\ne\ell$ satisfies $p\equiv N\pmod\ell$. (Then $q=N-p$ is divisible
by $\ell$, and the argument of Proposition~\ref{prop:twinblock} applies.)
A blocked class can exist only for $\ell=3$ and $3\nmid g$.
\end{proposition}

\begin{proof}
Such $g$-starters avoid the residues $0$ and $-g$ modulo $\ell$, which leaves
$\ell-2$ allowed classes (or $\ell-1$ if $\ell\mid g$). A class can be blocked
only if a single class is allowed, which forces $\ell-2=1$ and $\ell\nmid g$.
\end{proof}

For $\ell\ge5$, ruling out a blocked class needs at least two allowed
residue classes to contain starters. This is the expected situation, but proving it
is beyond current methods, since it asks for prime pairs in arithmetic
progressions. Proposition~\ref{prop:only3} should be read as saying that the
congruence argument has no analogue beyond $\ell=3$. It explains what we see:

\begin{itemize}
\item For $g\equiv 2\pmod 3$ with $3+g$ prime ($g=2,8,14,20,\dots$), the proof of
Proposition~\ref{prop:twinblock}(1) applies verbatim. Within the class
$N\equiv 2\pmod 6$, the single-gap version fails exactly when $N-3$ is composite.
For $g=8$, the full failure set up to $10^7$ coincides element by element with
that of $g=2$ ($1{,}334{,}283$ failures, all with $N\equiv 2 \pmod 6$).
\item For $3\mid g$ there is no blocked class. The single-gap versions with
$g=6$ and $g=12$ have no failure among even $8\le N\le 10^7$, and
Conjecture~\ref{conj:B} ($g=6$) is verified to $10^{12}$
(Section~\ref{sec:computation}).
\end{itemize}

\begin{remark}[blocking versus halving]\label{rem:halving}
For sexy starters the prime $3$ does not block, but it \emph{halves}. If
$N\equiv 1\pmod 3$ and $q>3$, then $p\equiv 2\pmod 3$; if $N\equiv 2\pmod 3$ and
$q>3$, then $p\equiv 1\pmod 3$; for $N\equiv0\pmod3$ both classes are usable.
Our data match this exactly: among $500{,}000$ minimal sexy witnesses up to
$10^{12}$, every $N\equiv1\pmod3$ used $p\equiv2$, and every $N\equiv2\pmod3$
used $p\equiv1$ except $N=8=5+3$. Table~\ref{tab:fuel} shows the resulting
cost per class.
\end{remark}

\begin{remark}\label{rem:infinitude}
If the set $S$ of starters (of any fixed kind) were finite, the Chinese
remainder theorem would give infinitely many even $N$ with $N-p$ divisible by
some prime $\ell_p$ and $N-p\neq \ell_p$ for every $p\in S$. Each such $N$ would be a
counterexample. Hence Conjecture~\ref{conj:A} implies infinitely many twin or
cousin pairs, and Conjecture~\ref{conj:B} implies infinitely many sexy pairs.
\end{remark}

\section{A density law}\label{sec:law}

Let $\pi(x;6,5)$ count primes $\le x$ that are $\equiv 5\pmod 6$.

\begin{theorem}\label{thm:law}
Let $B(X)$ be the number of even $N\le X$ with $N\equiv 2\pmod 6$ and $N-3$
composite (equivalently, by Proposition~\ref{prop:twinblock}, with no
twin-starter partition in that class). Then
\[
B(X)=\frac{X}{6}-\pi(X-3;6,5)+O(1)
=\frac{X}{6}-(1+o(1))\frac{X}{2\ln X},
\]
so the proportion of such $N$ among the even $N\le X$ is
$\tfrac13-(1+o(1))/\ln X$.
\end{theorem}

\begin{proof}
The map $N\mapsto N-3$ sends $\{N\le X:N\equiv2\pmod6\}$ bijectively onto
$\{m\le X-3: m\equiv5\pmod6\}$. The count follows, and the asymptotic is the prime
number theorem for arithmetic progressions,
$\pi(x;6,5)\sim x/(\varphi(6)\ln x)=x/(2\ln x)$~\cite{Davenport}.
\end{proof}

Locally, the density of the blocked set near $N$ is
$\tfrac13\,(1-3/\ln N)$: one third of all even numbers, minus the fraction
rescued by the $p=3$ route, which works with probability $3/\ln N$ because
$N-3$ lies in a class mod $6$ that contains half of all primes. Let $E_T$ be
the set of even $N$ with no twin-starter partition at all. Then $E_T$ contains
the blocked set, and the gearbox conjecture (Corollary~\ref{cor:gearbox}(a),(b) with
twin starters) says they are equal. Table~\ref{tab:law} compares the measured
density of $E_T$ with the band average of $\tfrac13(1-3/\ln N)$.

\begin{table}[ht]
\centering
\begin{tabular}{lrrr}
\toprule
Range of $N$ & Sample & Measured $|E_T|$ share & Predicted \\
\midrule
$[6,10^4]$ & all & $21.01\%$ & $20.91\%$\\
$(10^4,10^5]$ & all & $24.02\%$ & $24.02\%$\\
$(10^5,10^6]$ & all & $25.68\%$ & $25.67\%$\\
$(10^6,10^7]$ & all & $26.82\%$ & $26.82\%$\\
$(10^7,10^8]$ & all & $27.67\%$ & $27.67\%$\\
$[10^{12},10^{13}]$ & $9\times10^5$ & $29.935\%$ & $29.906\%$\\
\bottomrule
\end{tabular}
\caption{Fraction of even $N$ with no twin-starter partition, against the band
average of $\frac13(1-3/\ln N)$. The rows up to $10^8$ are exhaustive. The last row
uses the first $1000$ even $N$ of each of the $900$ stretches of length
$10^{10}$ (sampling standard error $\approx 0.05\%$). Within the class
$N\equiv2\pmod 6$ the measured share there is $89.81\%$, against
$1-3/\ln N$ averaged over the range, $89.72\%$.}
\label{tab:law}
\end{table}

The approach to $1/3$ is extremely slow: the shortfall is $\approx 1/\ln N$,
still $3.3\%$ at $N=10^{13}$. The local value $\tfrac13(1-3/\ln N)$ reaches
$30\%$ only at $N=e^{30}\approx 1.07\times10^{13}$, and $33\%$ only near $10^{130}$.
The limit itself will never be visible in computation; it is known only through
Theorem~\ref{thm:law}.

\section{Computation}\label{sec:computation}

\paragraph{Algorithm.}
The verifier is the Q-Hot method of~\cite{QHot,QHot2}, in the bitset form
of~\cite{QHot2}, with one change: its
\emph{anchor bitset}, which in~\cite{QHot2} marks every odd prime up to $P$, here
marks only the starters up to $P$ ($P=10^7$ or $2\cdot10^7$, see Table~\ref{tab:runs}; twin or cousin starters for
Conjecture~\ref{conj:A}, sexy starters for Conjecture~\ref{conj:B}). It is built
once by a sieve. Each thread keeps a \emph{ring} of up to $512$ large primes $q$
that recently completed a partition. The method rests on the observation of the
introduction: if $N=p+q$, then $N+2k=(p+2k)+q$ whenever $p+2k$ is again an
anchor, so one large prime serves a stretch of consecutive even numbers. For 64
consecutive even numbers $N_0,\dots,N_0+126$ and a cached $q<N_0$, the partners
$N_0-q+2i$ are 64 consecutive odd numbers, so a single 64-bit window of the
anchor bitset certifies all of them at once. No primality test is needed,
because $q$ is already proven prime. Numbers left uncovered go to a cold path,
which scans the anchors $p$ in ascending order, discards $q=N-p$ with a prime
factor up to $37$, and tests the rest with a deterministic seven-base
Miller--Rabin test valid for all $n<2^{64}$~\cite{Sinclair,FG}. The first $q$
found is proven prime and enters the ring, and the witness is the least anchor
partner of $N$. An $N$ with no anchor partner $p\le P$ is escalated to a search
over starters up to $10^8$ and logged as a miss. The range is cut into blocks of
$2^{b}$ even numbers ($b=20$ or $24$), each verified independently with a
freshly cleared ring, so results do not depend on the thread count. With all
odd primes as anchors, the same program verified the binary Goldbach conjecture
on $[4\cdot10^{18},4.003\cdot10^{18})$~\cite{QHot2}.

\paragraph{Integrity.}
The program writes checkpoints containing running totals, sample witnesses
(one fast-path and one cold-path witness per stretch of length $10^{10}$), and
every miss record, and it validates them on resume. A separate
log records the least starter for the first $w$ even $N$ of every stretch of
length $10^{10}$. Every build is checked by a bit-exact comparison of its
vectorised prefix against the scalar path, and by forced-miss runs whose miss
lists must agree between the profile-guided and plain builds. A provenance file
records compiler flags and SHA-256 hashes of source, binary and profile. Rebuilding
the Conjecture~\ref{conj:A} verifier from the final published source reproduced the
campaign binary bit for bit (identical SHA-256), so the published code is exactly the
code that produced the results. All
sampled witnesses, and all $9\times10^5$ and $5\times10^5$ logged minimal witnesses of
the campaigns on $[10^{12},10^{13}]$ (Conjecture~\ref{conj:A}) and $[8,10^{12}]$
(Conjecture~\ref{conj:B}), were rechecked by independent Python scripts (sympy
primality), and random subsamples were checked for minimality. For every logged
$N\equiv 2\pmod 3$ in the first of these, the audit also confirms
Proposition~\ref{prop:twinblock}: a twin-starter partner exists exactly when $N-3$
is prime.

\paragraph{Results.} Table~\ref{tab:runs} lists the campaigns. The two
ranges for Conjecture~\ref{conj:A} meet without gap or overlap.

\begin{table}[ht]
\centering
\begin{tabular}{llrrrr}
\toprule
Conjecture & Range & Even $N$ & Misses & Time (s) & Rate (M/s)\\
\midrule
\ref{conj:A} & $[6,10^{12}]$ & $499{,}999{,}999{,}998$ & $0$ & $994.7$ & $502.7$\\
\ref{conj:A} & $[10^{12}+2,10^{13}]$ & $4{,}500{,}000{,}000{,}000$ & $0$ & $8662.3$ & $519.5$\\
\ref{conj:B} & $[8,10^{12}]$ & $499{,}999{,}999{,}997$ & $0$ & $803.2$ & $622.5$\\
\midrule
\ref{conj:A} & $[4\cdot10^{18},4\cdot10^{18}+10^8]$ & $5\times10^7$ & $0$ & -- & --\\
\ref{conj:B} & $[4\cdot10^{18},4\cdot10^{18}+10^8]$ & $5\times10^7$ & $0$ & -- & --\\
\bottomrule
\end{tabular}
\caption{Verification campaigns. The anchor limit was $P=2\cdot10^7$ (block size $2^{24}$) for the
first row, $P=10^7$ (block size $2^{24}$) for the second, and $P=10^7$ (block size $2^{20}$) for the
others. The upper three rows ran
on a laptop with an Intel Core i5-1235U (10 cores, 12 threads) under WSL2,
compiled with GCC 13.3 (\texttt{-O3 -march=native}, AVX2 prefix); all three
rows used profile-guided builds (the Conjecture~\ref{conj:B} build with
\texttt{-DSEXY\_ANCHORS}). An earlier run of Conjecture~\ref{conj:B} with a plain,
non-profiled build gave identical counters (fast-path hits, cold-path candidates,
Miller--Rabin calls) to all reported digits, as the block design requires.
The last two rows are spot windows near the limit of published Goldbach
verification~\cite{OHP}.}
\label{tab:runs}
\end{table}

\section{How small the starter is}\label{sec:fuel}

\begin{table}[ht]
\centering
\begin{tabular}{lrrrr}
\toprule
& \multicolumn{2}{c}{Conjecture~\ref{conj:A}} & \multicolumn{2}{c}{Conjecture~\ref{conj:B}}\\
\cmidrule(lr){2-3}\cmidrule(lr){4-5}
$N\bmod 3$ & mean & max & mean & max \\
\midrule
$0$ & $105$ & $1{,}783$ & $110$ & $2{,}441$\\
$1$ & $311$ & $6{,}299$ & $268$ & $6{,}257$\\
$2$ & $261$ & $6{,}007$ & $332$ & $8{,}641$\\
\bottomrule
\end{tabular}
\caption{Least starter partner $p$ for $2\times10^4$ consecutive even $N$ from
$4\cdot10^{18}$, by residue class. Under Conjecture~\ref{conj:A}, each blocked
class is served by one dedicated gap; under Conjecture~\ref{conj:B}, classes $1$
and $2$ each lose half of the starters (Remark~\ref{rem:halving}). The overall
means are similar: $226$ (A) and $237$ (B).}
\label{tab:fuel}
\end{table}

The least starter stays tiny as $N$ grows. Over the sampled windows, its
maximum was $4{,}729$ up to $10^{12}$ and $6{,}197$ on $[10^{12},10^{13}]$ for
Conjecture~\ref{conj:A}, and $4{,}357$ up to $10^{12}$ for Conjecture~\ref{conj:B}.
Near $4\cdot10^{18}$ it was at most $6{,}299$ and $8{,}641$ respectively.
These are maxima over sampled windows only, and they understate the exhaustive maxima,
because the extreme $N$ are rare. The class $N\equiv 1\pmod 3$ can be compared exactly with
the exhaustive records of~\cite{Pandey}.

\begin{proposition}\label{prop:samewitness}
Let $N\equiv 1\pmod 3$ be even. Among partitions $N=p+q$ with $q>3$ prime, $p$ is a twin or
cousin starter if and only if $p$ is a lower twin member, if and only if $p$ is a twin member.
Hence the least twin-or-cousin starter partner of $N$ equals the least twin-member partner
$t_{\min}(N)$ of~\cite{Pandey}, unless one of them equals $N-3$.
\end{proposition}

\begin{proof}
A cousin starter $p>3$ has $p\equiv 1\pmod 3$, which would give $3\mid q$; and $3$ is a twin
starter. An upper twin member $p$ with $q>3$ is $5$ by~\cite[Prop.~3]{Pandey}, and $5$ is
also a lower member. So all three sets of admissible $p$ coincide.
\end{proof}

All $21$ of the record values of $t_{\min}$ in~\cite{Pandey} with $N\equiv1\pmod 3$ agree
with the least twin-or-cousin starter computed independently by us (script
\texttt{crosscheck\_pandey\_records.py}); those records were confirmed to $10^{13}$ by
Oliveira e Silva with independent code~\cite{Pandey}. In particular the exhaustive maximum
of the least starter in Conjecture~\ref{conj:A} is at least $20{,}747$ below $10^{13}$
(at $N=6{,}527{,}240{,}154{,}856$) and at least $21{,}599$ below $10^{14}$. In the class
$N\equiv2\pmod3$ the two problems differ sharply: at Pandey's records in that class, which
need upper twin members, our least starter (a cousin starter) is between $37$ and $937$.

The two starter sets are of comparable size. Below $10^7$ there are
$117{,}601$ twin-or-cousin starters and $117{,}207$ sexy starters. This
agrees with the Hardy--Littlewood prediction~\cite{HL}: sexy pairs are
asymptotically twice as numerous as twin pairs, and twin and cousin pairs are
equally numerous and (apart from $p=3$) disjoint as starter sets.

In the data for Conjecture~\ref{conj:B} up to $10^{12}$, the class means are
unequal: $51$, $122$ and $167$ for $N\equiv 0,1,2\pmod 3$. This is because the smallest
sexy starters lean towards $p\equiv 2\pmod 3$ ($5,11,17,23$), which favours the
class $N\equiv1$. Near $4\cdot10^{18}$ the asymmetry between classes $1$ and
$2$ is smaller. We expect it to be a small-number effect, but we have not
quantified it.

\section{Open questions}
\begin{enumerate}
\item Prove Conjecture~\ref{conj:A} or~\ref{conj:B} for almost all even $N$,
perhaps conditionally on a Hardy--Littlewood-type hypothesis, in the spirit of
the conditional density result of~\cite{Pandey}.
\item Give an asymptotic for the number of starter partitions of $N$ (a
singular series for the restricted problem), and explain the class-dependent
costs of Table~\ref{tab:fuel}.
\item Is the set of exceptions to the single-gap version finite for every
$g\equiv0\pmod 6$? (For $g=6$ and $g=12$ the only exception up to $10^7$ is
$N=6$.)
\item Extend the verification of Conjecture~\ref{conj:A} beyond $10^{13}$ and of
Conjecture~\ref{conj:B} beyond $10^{12}$.
\end{enumerate}

\section*{Data and code availability}
The verifier source, build script, analysis scripts, checkpoint files, minimal-witness
logs and provenance files are available at \url{https://github.com/psaraj12/goldbach-engine}, in the directory
\texttt{goldbach-starters/}.

\section*{Acknowledgements}
% EDIT as you see fit; disclosure of AI assistance is recommended.
Code review, analysis scripts and drafting of this manuscript were assisted by an
AI system (Claude, Anthropic). The author designed the conjectures, ran all
computations and checked every result.

\begin{thebibliography}{99}
\bibitem{A007534} OEIS Foundation Inc., Sequence A007534, \emph{Positive even
numbers that are not the sum of a pair of twin primes},
\url{https://oeis.org/A007534}.
\bibitem{Davenport} H. Davenport, \emph{Multiplicative Number Theory}, 3rd ed.,
Graduate Texts in Mathematics 74, Springer, 2000.
\bibitem{Dubner} H. Dubner, Twin prime conjectures, \emph{J. Recreational Math.}
\textbf{30} (1999--2000), no.~3.
\bibitem{Gomez} J. Gómez, On primorial numbers, arXiv:2301.02770.
\bibitem{HL} G. H. Hardy and J. E. Littlewood, Some problems of `Partitio
numerorum' III: On the expression of a number as a sum of primes, \emph{Acta
Math.} \textbf{44} (1923), 1--70.
\bibitem{OHP} T. Oliveira e Silva, S. Herzog and S. Pardi, Empirical
verification of the even Goldbach conjecture and computation of prime gaps up to
$4\cdot10^{18}$, \emph{Math. Comp.} \textbf{83} (2014), 2033--2060.
\bibitem{Pandey} T. Pandey, An intermediate conjecture between Goldbach and
Dubner: every even number is the sum of a prime and a twin prime,
arXiv:2608.02381 (2026).
\bibitem{Sahu} R. Sahu, On extending Goldbach's conjecture: the patterns of
isolated and twin primes into the even sums, \emph{Int. Res. J. Eng. Technol.}
\textbf{11} (2024), no.~10, 62--65.
\bibitem{Sinclair} J. Sinclair (2011), the seven-base set
$\{2, 325, 9375, 28178, 450775, 9780504, 1795265022\}$, deterministic for all
$n<2^{64}$; recorded in W. Izykowski, \emph{Deterministic variants of the
Miller--Rabin primality test}, \url{https://miller-rabin.appspot.com}.
\bibitem{FG} J. Feitsma and W. Galway, list of base-2 strong pseudoprimes below
$2^{64}$, \url{http://www.cecm.sfu.ca/Pseudoprimes/index-2-to-64.html}.
\bibitem{QHot} S. A. R. Parthibanathan, Q-Hot Cache: a high-throughput method for
empirical verification of the even Goldbach conjecture, Zenodo preprint, April 2026,
doi:\href{https://doi.org/10.5281/zenodo.19884541}{10.5281/zenodo.19884541}.
\bibitem{QHot2} S. A. R. Parthibanathan, Q-Hot Cache: verifying the even Goldbach
conjecture to $4.003\cdot10^{18}$, submitted to \emph{Experimental Mathematics} (2026); campaign data
doi:\href{https://doi.org/10.5281/zenodo.23082138}{10.5281/zenodo.23082138}, production audit
doi:\href{https://doi.org/10.5281/zenodo.23278480}{10.5281/zenodo.23278480}. % CHECK: replace with arXiv/journal ref once posted
\bibitem{Sun} Z.-W. Sun, Conjectures on representations involving primes, in:
M. Nathanson (ed.), \emph{Combinatorial and Additive Number Theory II},
Springer Proc. Math. Stat. \textbf{220}, Springer, New York, 2017, pp.~279--310;
arXiv:1211.1588.
\bibitem{Zwillinger} D. Zwillinger, A Goldbach conjecture using twin primes,
\emph{Math. Comp.} \textbf{33} (1979), 1071,
doi:\href{https://doi.org/10.1090/S0025-5718-1979-0528060-5}{10.1090/S0025-5718-1979-0528060-5}.
\end{thebibliography}

\end{document}
